\chapter{Conclusion}

\section{Summary}
%szb: nem kell setion, csak a chapter

This thesis investigated the effectiveness of large language models for
Java-to-C++ code translation. Two code-specific model families, StarCoder2 and
Qwen2.5-Coder, were evaluated in both their base and fine-tuned forms using
BLEU score and execution-based evaluation with CodeEditorBench.

The results showed that fine-tuning improved the performance of both models in
terms of textual similarity and functional correctness. Among the evaluated %szb: ne azon legyen a hangsúly, hogy melyik teljesített jobban, hanem az elsődlegesen vizsgált aspektusokon
models, the fine-tuned Qwen2.5-Coder model achieved the strongest overall
performance, producing the highest BLEU score and the largest number of accepted
solutions.

The experiments also demonstrated that similarity-based metrics such as BLEU are
not sufficient for evaluating code translation systems on their own. Several
generated programs achieved relatively high BLEU scores while still failing to
compile or execute correctly. This supports earlier research showing that
execution-based evaluation provides a more reliable measure of functional
correctness than token-level similarity metrics.

Although the fine-tuned models performed substantially better than their base
versions, the number of remaining compilation errors and wrong answers indicates
that fully automated code translation is not yet reliable enough for practical
use without human supervision. Current models can already support developers by
generating initial translations, but human review is still necessary to ensure
correctness, maintainability, and compatibility with the target language.